\section{Cote 136: removing the sum over primes from the finite gluing bound}
\label{sec:136arithmetic}

\paragraph{Question and scope.}
This is a refinement of ALPOGE\_BRIDGES AQ-136-SEMILOCAL, concerning
the uniform weight bounds of physical pp.~117--121, footers
116--120, scans 117, 119 and 121.  It does not reprove the already
accepted semilocal terminality theorem.  It changes its construction
of a common source: one rational comparison theory suffices, and the
term $\sum_{p\in P}N_p$ from (136.S2) disappears.  The remaining
unbounded local scales $N_p$ still prevent this construction from
answering conjecture (a) over $\mathbf Z$.

\paragraph{ARCH-136-ARITH-BOUND.}
Let $P=\{p_1,\ldots,p_s\}$ be finite and nonempty, let $k$ be an
$R_P=\mathbf Z_{(P)}$-algebra, and set
$k_i=k\otimes_{R_P}\mathbf Z_{(p_i)}$, $K=k\otimes_{R_P}\mathbf Q$.
Fix a finite nonempty family $C_v\in\mathcal A_k$ and choose $N\geq1$
such that every $(C_v)_*^{w,0}$ is acyclic whenever $w\geq N$.
Choose $N_i=p_i^{a_i}\geq N$ and write $M=\max_iN_i$.
There is an admissible arithmetic comparison object $A$ with:
\begin{enumerate}
\item horizontal acyclicity in form degree $q$ for $w\geq(q+1)M$;
\item coefficient vertical cohomology through degree $d$ for
      $w\geq N(d+2)$, uniformly in every simplex;
\item strict maps $A\to\mathcal E_J(C_v)$, where each
      $C_v\to\mathcal E_J(C_v)$ is a natural arrow to an admissible
      arithmetic replacement using a fixed two-edge simplicial interval
      $J$.  The admissibility cutoffs of this replacement depend on
      those of $C_v$, and can be chosen independently of $s$.
\end{enumerate}
Consequently actual strict common admissible sources of the $C_v$
can be obtained by path pullbacks with horizontal cutoffs depending
on $M$ and the original target cutoffs, but with no sum over $i$.
Their vertical cutoffs can be chosen independently of the numbers
$N_i$ and $s$, once $N$ and the finite target family are fixed.

\begin{proof}
If $K=0$, all rational objects and maps below are zero; the same
exact arithmetic-sequence argument applies.  Thus this case does
not require dividing by a nonzero element of $K$.

Put $A_i=B_{N_i}^{\mathrm{PD}}(k_i)$, using AQ-136-PLOCAL, and
put $F=B_N(K)$, the rational source whose coordinate relation is
$\sum_jx_j=T^{[N]}$.  Rescaling by $N!$ identifies it with the
version whose relation is $t^N$.  These objects have exact vertical
cohomology equal to their coefficient algebra in every weight.
Their horizontal cutoffs are respectively $(q+1)N_i$ and $(q+1)N$.

Since $N_i\geq N$, the weight-$N_i$ zero-form row of $F$ is
acyclic.  The coordinate lifting argument of AQ-136-PLOCAL supplies
a strict rational arrow
\[
                         j_i:(A_i)_K\longrightarrow F.
\]
Indeed the coefficient $T^{[N_i]}$ has a normalized degree-one
primitive in that row, which defines a linear map from the standard
coordinate disk; the free differential PD presentation extends it
to the displayed arrow.  This does not require $N_i$ to be a
multiple of $N$.

Similarly choose strict $a_{i,v}:A_i\to(C_v)_i$ and
$b_v:F\to(C_v)_K$.  The two rational maps
$(a_{i,v})_K$ and $b_vj_i$ out of $(A_i)_K$ are simplicially
homotopic relative to the coefficients.  For precision, their
coordinate difference in weight $N_i$ vanishes on the constant
submodule of the standard coordinate disk $L$.  It factors through
$Q=L/K$, with normalized complex $K[1]$.  Its image cycle is a
boundary because the target weight-$N_i$ zero-form row is acyclic.
Dold--Kan gives a linear nullhomotopy; adding the first coordinate
map and extending by the free PD presentation gives a strict
\[
 H_{i,v}:(A_i)_K\longrightarrow P_{\Delta[1]}(C_v)_K
\]
with the required endpoints and constant coefficient paths.

Define the arithmetic object by strict algebraic limits:
\[
 A=\left\{((a_i),a_0,(\omega_i)):
 \begin{array}{l}
 a_i\in A_i,\ a_0\in F,\ \omega_i\in P_{\Delta[1]}F,\\
 e_0\omega_i=j_i(a_i|_K),\quad e_1\omega_i=a_0
 \end{array}\right\}.
                                                        \tag{136.A1}
\]
The coefficient map is diagonal localization with constant paths.
It is a strict simplicial differential PD $S_k$-algebra.
In form degree $q$, endpoint surjectivity for the $F$ row holds
for $w\geq(q+1)N$, by ARCH-136-PATHS.  Hence in these weights
there is an exact sequence of underlying simplicial modules
\[
 0\to A^{w,q}\to
 \left(\bigoplus_iA_i^{w,q}\right)\oplus F^{w,q}
       \oplus(P_{\Delta[1]}F^{w,q})^{\oplus s}
 \to(F^{w,q})^{\oplus2s}\to0.                            \tag{136.A2}
\]
All rows after the first are acyclic for $w\geq(q+1)M$.
The normalized long exact sequence proves assertion 1.

To prove assertion 2 through degree $d$, it is enough to take
$w\geq N(d+2)$.  The comparison proof AQ-136-COMPARISON for $F$
uses vertical exactness in every weight and horizontal acyclicity
through form degree $d+1$.  Those hypotheses hold at this weight,
so $P_{\Delta[1]}F$ has coefficient cohomology through degree $d$,
uniformly in the simplex.  Sequence (136.A2) is degreewise exact
through degree $d+1$.  Replace its last term, if necessary, by the
image in all degrees; through degree $d$ that image has the same
cohomology.  On cohomology in degree zero, eliminate the path
constant for each $i$.  Kernel and cokernel are then exactly those
of the arithmetic sequence
\[
 0\to S_k^w\to
       \left(\bigoplus_iS_{k_i}^w\right)\oplus S_K^w
       \to(S_K^w)^{\oplus s}\to0,
\]
which is AQ-136-SEMILOCAL, (136.S1).  It is exact for arbitrary
$k$-modules.  All positive cohomology of the other terms through
$d$ vanishes.  The long exact sequence gives the claimed result.
Notice that this vertical argument uses no bound involving $N_i$.

For an arbitrary admissible $C$, define $\mathcal E_J(C)$ by the
same arithmetic limit with local values in $C_i$, central value
in $C_K$, and paths in $P_J(C_K)$ joining the rationalized local
value to the central value.  The natural $C\to\mathcal E_J(C)$
uses its localizations and constant paths.  For a fixed form degree
$q$, every constituent has the same eventual horizontal acyclicity
bound as $C^{*,q}$, and endpoint surjectivity holds at that bound.
For a fixed vertical range, AQ-136-COMPARISON supplies one uniform
bound for all the path values; take also the finitely many required
row bounds through $d+1$.  The same arithmetic-sequence argument
proves admissibility.  Finite direct sums, however many, do not
enlarge these numerical bounds.  This proves the claimed uniformity
in $s$ for $\mathcal E_J$.

Take $J$ to be two intervals glued end to end.  The map
$A\to\mathcal E_J(C_v)$ sends local values by $a_{i,v}$, its
central value by $b_v$, and its $i$th path to the concatenation
\[
                   H_{i,v}(a_i|_K)\ \text{ followed by }\ b_v(\omega_i).
\]
The common endpoint is exactly $b_vj_i(a_i|_K)$, so the sections
glue.  Compatible sections on a glued simplicial set form a
pullback; consequently this concatenation is a strict algebraic
map, preserving products, differential and divided powers.  Its
endpoints are precisely the specified local and central values.
This proves assertion 3.

Finally apply AQ-136-COSPAN to $A\to\mathcal E_J(C_1)\leftarrow C_1$,
then to the resulting source and $C_2$ over $\mathcal E_J(C_2)$,
and continue through the finite family.  The projections remain
strict arrows to every previous target.  In the horizontal argument
for each pullback only maxima of the form-$q$ row bounds occur;
there is no sum over the primes.  Its vertical argument takes a
finite maximum of vertical and row bounds through $d+1$ for the
middle object and the vertical bounds for its two endpoints.
The middle object is $\mathcal E_J(C_v)$, whose bounds do not
involve the $N_i$; the vertical bounds for $A$ already have the
same independence.  Induction through the fixed finite target
family therefore gives the asserted vertical independence for
the final strict common source.\hfill$\square$
\end{proof}

\paragraph{What remains over $\mathbf Z$.}
For a fixed integral family, $N$ can be chosen before the finite
prime set.  The theorem proves that the rational tensor term
$\sum_iN_i$ is not intrinsic to arithmetic comparison.  The local
term $\max_iN_i$ remains, and it cannot be suppressed while retaining
the same prescribed local models: for $p>N$ their scale is at least
$p$, so in weights $w<p$ their zero-form row is the nonzero constant
coefficient module.  An integral admissible source cannot map
after localization to all these objects, because its high-weight
coefficient is a horizontal boundary, whose image cannot be the
nonzero constant coefficient class.  This exact obstruction was
already recorded in checkpoint~09, \texttt{root\_136\_routes\_audit.md},
``Infinite arithmetic gluing test''; it is used, not claimed again
as a new result.

Thus the remaining arithmetic construction problem is to replace
the local sources themselves by models with a common horizontal
weight bound, or to use a construction without arrows to those
prescribed local sources.  Merely choosing a smaller rational
comparison object is insufficient.  This distinction concerns
uniform resolutions by generalized divided-power de Rham theories
over $\operatorname{Spec}\mathbf Z$, not the validity of the
fixed-prime crystalline Poincare lemma.
